% Generated draft: requires mathematical and source review.
\documentclass[11pt]{article}
\usepackage[margin=25mm]{geometry}
\usepackage{fontspec}
\setmainfont{Latin Modern Roman}
\usepackage{amsmath,amssymb,mathtools,unicode-math}
\setmathfont{Latin Modern Math}
\usepackage{enumitem,needspace,xurl,hyperref}
\setlength{\parindent}{0pt}
\setlength{\parskip}{5pt}
\setlength{\emergencystretch}{3em}
\newcommand{\N}{\mathbb{N}}
\newcommand{\Z}{\mathbb{Z}}
\newcommand{\Q}{\mathbb{Q}}
\newcommand{\R}{\mathbb{R}}
\newcommand{\C}{\mathbb{C}}
\newcommand{\F}{\mathbb{F}}
\newcommand{\A}{\mathbb{A}}
\newcommand{\PP}{\mathbb{P}}
\newcommand{\E}{\mathbb{E}}
\newcommand{\eps}{\varepsilon}
\newcommand{\dd}{\,\mathrm{d}}
\begin{document}
\section*{UnsolvedMath ID: 20000312}
\textbf{Source problem code:} catalogue entry
\subsection{Definitions and mathematical statement}
In the intended characteristic-zero setting, let $x\in X$ be a point on a normal projective $n$-fold with $K_X$ $\Q$-Cartier near $x$. Define its local canonical index by
\[\operatorname{ind}_x(K_X)=\min\{r\in\Z_{>0}:rK_X\text{ is Cartier near }x\}.\]
For a divisor $E$ over $X$, write $A_X(E)=1+\operatorname{coeff}_E(K_Y-f^*K_X)$; the exact-center invariant is $\operatorname{mld}_x(X,0)=\inf_{c_X(E)=x}A_X(E)$. Canonical and terminal singularities have respectively $A_X(E)\ge1$ and $A_X(E)>1$ for exceptional divisors. A point that is itself a crepant center is the center of a divisor of log discrepancy one.

\textbf{Original questions (Shokurov).} For fixed $n$, is the local index bounded for projective varieties with an ``exact canonical'' singularity of dimension $n$? In particular, what about $\dim X=3$? More generally, for fixed dimension $n$ and a fixed value $a$ of the exact-center minimal log discrepancy, is there a bound $R(n,a)$ such that
\[\operatorname{mld}_x(X,0)=a\quad\Longrightarrow\quad\operatorname{ind}_x(K_X)\le R(n,a)\]
for every point in the intended class?

The phrase ``exact canonical'' is undefined in the original. The research aid says it is a typo for ``strictly canonical'' in an earlier AIM draft, and interprets the meaningful local case as a canonical point that is itself a crepant center, equivalently exact-center log mld $1$. Being nonterminal somewhere in a neighborhood does not by itself specify that center. The earlier wording and intended equivalence have not been independently source-verified. The characteristic, precise singularity class for the general fixed-mld question, and point-versus-subvariety center conventions remain unstated.

\textbf{Attached unreviewed research-aid account.} The aid attributes to Kawakita the sharp index bound $6$ in dimension three for the indicated log-mld-one canonical interpretation, and fixed-mld index boundedness for canonical threefolds. It reports the general-dimensional conjecture as open and the quotient case as known, without exact hypotheses or citations.

It additionally gives a varying-mld family: for every $n\ge3$ and every integer $r>n-2$ with $\gcd(r,n-2)=1$, consider, over the intended characteristic-zero field containing the $r$th roots of unity,
\[X_{n,r}=\PP^n/\mu_r(0,1^{n-1},r-(n-2)).\]
Here the weights specify the diagonal action on homogeneous coordinates; the chart corresponding to the weight-zero coordinate contains the germ at the origin of
\[U_{n,r}=\A^n/\mu_r(1^{n-1},r-(n-2)).\]
The report claims this germ is isolated terminal and
\[\operatorname{mld}_0(U_{n,r})=1+\frac1r,\qquad\operatorname{ind}_0(K_{U_{n,r}})=r,\]
with smooth index-one cover $\A^n$ of degree $r$. Its claimed calibration is
\[\deg(\text{index-one cover})=\operatorname{ind}_0(K_{U_{n,r}})=\frac1{\operatorname{mld}_0(U_{n,r})-1}=r.\]
Consequently any fixed-mld bounding function would satisfy $R(n,1+1/r)\ge r$ along these admissible parameters. The mld varies with $r$; this family concerns dependence of a bound on $a$, rather than being a counterexample with one fixed mld value or with log mld exactly one. All family and literature assertions are supplied research-aid claims requiring review.
\subsection{Short English statement}
Shokurov asks for local canonical-index bounds at a fixed dimension, first at an ambiguously named canonical singularity and then at fixed mld. The aid interprets the first as an exact-center mld-one question, reports threefold results, and supplies terminal quotient germs whose indices grow while their mlds approach one from above.
\subsection{Sources}
\begin{itemize}
\item Local UnsolvedMath registry, permanent ID 20000312, source code catalogue entry; \texttt{aim-algebraic-geometry-notes.json}, notes[311], tag \texttt{problem}.
\item AIM workshop \emph{Numerical invariants of singularities and higher-dimensional algebraic varieties}, Problem 1, attributed to Shokurov: \url{https://aimath.org/WWN/singularvariety/problems_final.pdf}. The extracted text says ``exact canonical''; the earlier ``strictly canonical'' draft is mentioned only in the research aid. No PDF was consulted.
\item Attached machine-generated research attempt 1 names Kawakita without bibliographic details, and gives no citation for the quotient-case status claim or earlier draft. Result \texttt{partial\_result}, status \texttt{partially\_solved}, rigor label \texttt{proved}, novelty confidence low. It is not peer reviewed and requires independent verification before citation. The original literature field is empty.
\end{itemize}
\end{document}
